\section{Grobe Flow-Analyse}
\label{sec:flow}

Die Aufgabe dieses Teils ist eine andere als in einer Analyse mit
vorgegebenen Handlungsoptionen. Es gibt hier keine Alternativen zu
vergleichen: Die Entscheidung ist gefallen, die \"Ubernahme bezahlt und die
Verschuldung eingegangen. Zu beantworten sind deshalb zwei Fragen. Erstens,
ob der Konzern die aufgenommenen Mittel aus eigener Kraft zur\"uckf\"uhren
kann und wie lange das dauert. Zweitens, welche Rendite ein Anleger
erwarten darf, der die Aktie zum heutigen Kurs kauft.

\subsection*{Annahmen der Flow-Analyse}
\addcontentsline{toc}{subsection}{Annahmen der Flow-Analyse}

Alle folgenden Annahmen stehen vor der Rechnung, die sie verbraucht. Sie
sind in \texttt{scripts/flow.py} an genau einer Stelle hinterlegt; die
Tabellen dieses Teils sind von diesem Skript erzeugt und nicht abgetippt.

\annahme{\textbf{Ausgangsbestand und Ausgangsverschuldung.} Ausgangspunkt
sind die eigenen Zahlungsmittel von \MioUSD{\Zahlungsmittel} und die
Finanzverbindlichkeiten von \MioUSD{\Gesamtverschuldung} zum
31.\,Dezember 2025. Die Treuhandmittel von \MioUSD{\TreuhandMittel} bleiben
au{\ss}er Ansatz, weil sie Kunden und Versicherern geh\"oren und zur Tilgung
nicht zur Verf\"ugung stehen.}

\annahme{\textbf{Umsatzbasis.} Fortgeschrieben wird nicht der berichtete
Umsatz von \MioUSD{\UmsatzGesamt}, sondern der Pro-forma-Umsatz von
\MioUSD{\ProFormaUmsatz}. Die \"Ubernahme wurde erst zum 1.\,August 2025
konsolidiert; der berichtete Umsatz enth\"alt sie nur f\"unf Monate und
w\"urde die Ertragskraft, aus der die Schulden zu bedienen sind, um rund ein
Sechstel untertreiben.}

\annahme{\textbf{Organisches Wachstum.} Die drei Szenarien sind keine
Sch\"atzungen des Verfassers, sondern drei vom Unternehmen selbst
ver\"offentlichte Wachstumsraten, jeweils konstant fortgeschrieben:
\Pct{\HJOrganisch} (erstes Halbjahr 2026), \Pct{\OrganischWachstum}
(Gesch\"aftsjahr 2025) und \Pct{\OrganischWachstumVJ} (Gesch\"aftsjahr
2024). Das mittlere Szenario tr\"agt die Haupttabellen. Damit muss keine
Wachstumsannahme erfunden werden, und die Spannweite der Rechnung ist die
Spannweite der eigenen j\"ungeren Erfahrung des Unternehmens.}

\annahme{\textbf{Marge.} Die bereinigte EBITDAC-Marge von \Pct{\FlowMarge}
des Jahres 2025 wird unver\"andert fortgeschrieben. Weder ein
Synergieeffekt noch eine Verw\"asserung durch die \"Ubernahme wird
unterstellt; beides w\"are nicht belegbar. Im ersten Halbjahr 2026 lag die
Marge bei \Pct{\HJEBITDACberMarge} und damit nahe an diesem Wert.}

\annahme{\textbf{Zins und Tilgung.} Beide folgen dem im Jahresbericht
ver\"offentlichten F\"alligkeitsprofil: \MioUSD{\ZinsverpflichtungJahrEins}
Zinsen im ersten Jahr, \MioUSD{\ZinsverpflichtungZweiDrei} f\"ur die Jahre
zwei und drei und \MioUSD{\ZinsverpflichtungVierFuenf} f\"ur die Jahre vier
und f\"unf, bei planm\"a{\ss}igen Tilgungen von
\MioUSD{\TilgungJahrEins}, \MioUSD{\TilgungJahrZweiDrei} und
\MioUSD{\TilgungJahrVierFuenf}.\quelleGB{2025}{46}\quelleMerken{s12gb} Der
Zinsaufwand sinkt in der Rechnung nicht mit einer vorzeitigen Tilgung. Die
Rechnung untertreibt die Entschuldung dadurch, statt sie zu
besch\"onigen.}

\annahme{\textbf{Abschreibungen und Sachinvestitionen.} Beide werden auf die
Jahresrate des ersten Halbjahres 2026 hochgerechnet:
\MioUSD{\FlowAbschreibung} Abschreibungen und \MioUSD{\FlowSachinvest}
Sachinvestitionen.\quelleQM{Form 10-Q Q2 2026}{29}\quelleMerken{s12qmformqq} Die Werte des Jahres 2025 enthalten
die \"Ubernahme nur f\"unf Monate und w\"aren zu niedrig. Die Abschreibung
bleibt \"uber den Horizont konstant, obwohl sie mit fortschreitender
Amortisation der Kundenbeziehungen sinken wird; das mindert die
ausgewiesene Steuerlast und ist damit die vorsichtigere Annahme.}

\annahme{\textbf{Dividende und Erwerbspreisnachzahlungen.} Die Dividende
folgt dem Quartalsbeschluss vom 21.\,Januar 2026, also
\USDje{\FlowDividendeJeAktie} im Jahr auf \num{\HJAktienUmlauf} Millionen
Aktien und damit \MioUSD{\FlowDividende}; eine weitere Erh\"ohung wird nicht
unterstellt, weil keine beschlossen ist. Die Erwerbspreisnachzahlungen
folgen dem ver\"offentlichten Profil von \MioUSD{\EarnoutMaximum}
Gesamtvolumen.\quelleErneut{s12gb}}

\annahme{\textbf{Keine weiteren Zuk\"aufe.} Dies ist die weitreichendste
Annahme der gesamten Rechnung und zugleich die unrealistischste: Das
Unternehmen hat allein 2025 \num{\AkquisitionenAnzahl} Gesellschaften
erworben, und Zuk\"aufe sind sein Gesch\"aftsmodell. Die Rechnung ist
deshalb keine Prognose, sondern eine Kapazit\"atsrechnung. Sie beantwortet
die Frage, wie schnell der Konzern die \"Ubernahme abtragen \emph{k\"onnte},
wenn er alle freien Mittel daf\"ur einsetzt, und macht damit den Spielraum
sichtbar, der stattdessen f\"ur weitere Zuk\"aufe zur Verf\"ugung steht.}

\annahme{\textbf{Steuerliche Betrachtung.} Anders als eine reine
EBITDA-Rechnung ist diese Analyse eine Nachsteuerbetrachtung: Zinsen und
Abschreibungen werden vom EBITDAC abgezogen und auf das verbleibende
Ergebnis die Konzernsteuerquote 2025 von \Pct{\FlowSteuerquote} angewendet.
Die ausgewiesenen internen Zinsf\"u{\ss}e sind deshalb Nachsteuerrenditen
auf Unternehmensebene; Steuern des Anlegers selbst bleiben unber\"ucksichtigt.}

\annahme{\textbf{Renditehorizont.} Der interne Zinsfu{\ss} wird \"uber
f\"unf, zehn und f\"unfzehn Jahre ausgewiesen und nie ohne seinen Horizont.
Er wird durch Bisektion bestimmt, nicht gesch\"atzt.}

\subsection{Entschuldungspfad des Konzerns}

Tabelle~\ref{tab:flowpfad} schreibt den Konzern unter dem mittleren
Szenario bis 2030 fort. Der nach Zins, Steuer, Sachinvestition, Dividende
und Erwerbspreisnachzahlung verbleibende Betrag geht dabei vollst\"andig in
die Schuldentilgung.

\begin{table}[htbp]
\centering
\caption{Entschuldungspfad, organisches Wachstum \Pct{\OrganischWachstum}
je Jahr. Betr\"age in Mio.~USD, Verschuldungsgrad als Vielfaches des
bereinigten EBITDAC.}
\label{tab:flowpfad}
\small
\begin{tabular}{lrrrrrrrr}
\toprule
Jahr & EBITDAC & Zinsen & Steuern & Earn-out & Tilgung & Freier CF & Nettoschuld & Grad\\
\midrule
\tabellenkoerper{data/flow_pfad}
\bottomrule
\end{tabular}
\end{table}

\FloatBarrier

Im ersten Jahr bleiben nach allen Auszahlungen \MioUSD{\FlowFreiJahrEins}
\"ubrig; der Betrag steigt bis 2030 auf ein Mehrfaches, weil die Zinslast
dem F\"alligkeitsprofil folgend sinkt und keine
Erwerbspreisnachzahlungen mehr anfallen. Der Verschuldungsgrad f\"allt von
\Faktor{\FlowGradStart} auf \Faktor{\FlowGradEnde}.

\subsection{Empfindlichkeit gegen\"uber dem organischen Wachstum}

Der Verschuldungsgrad ruht auf einer Wachstumsannahme, und diese Annahme
ist der strittige Punkt des gesamten Falls.
Tabelle~\ref{tab:flowszenarien} zeigt ihn deshalb f\"ur alle drei
Szenarien.

\begin{table}[htbp]
\centering
\caption{Nettoverschuldung je Einheit bereinigtes EBITDAC am Jahresende, je
Szenario. Die Rechnung endet bei null; dar\"uber hinausgehende Mittel
stehen f\"ur Zuk\"aufe oder R\"uckk\"aufe zur Verf\"ugung.}
\label{tab:flowszenarien}
\begin{tabular}{lrrrrrr}
\toprule
Szenario & Wachstum in \% & 2026 & 2027 & 2028 & 2029 & 2030\\
\midrule
\tabellenkoerper{data/flow_szenarien}
\bottomrule
\end{tabular}
\end{table}

\FloatBarrier

\bk{Die Rangfolge ist \"uber alle drei Szenarien stabil, und das Niveau
ver\"andert sich weniger, als die Spannweite der Wachstumsraten vermuten
l\"asst. Selbst im ung\"unstigsten Fall, der das negative organische
Wachstum des ersten Halbjahres 2026 auf f\"unf Jahre fortschreibt, sinkt der
Verschuldungsgrad von \Faktor{\FlowGradStart} auf
\Faktor{\FlowGradSzenarioA}. Der Grund liegt in der Struktur des
Gesch\"afts: Ein Makler bindet kaum Kapital, seine Sachinvestitionen von
\MioUSD{\FlowSachinvest} betragen weniger als drei Prozent des EBITDAC, und
sein Cash-Flow ist deshalb fast vollst\"andig frei verf\"ugbar. Die
Verschuldung ist damit tragbar; sie ist nicht das Risiko dieses Falls.}

\subsection{Rendite des Anlegers}

Die zweite Rechnung wechselt die Sicht. Ein Anleger kauft die Aktie zum
Schlusskurs vom 27.\,August 2026 von \USDje{\FlowEinstiegskurs}, erh\"alt
in jedem Jahr den aussch\"uttbaren Cash-Flow je Aktie und verkauft am Ende
des Horizonts.

\annahme{\textbf{Abgrenzung gegen den Entschuldungspfad.} In dieser Rechnung
wird nur die \emph{planm\"a{\ss}ige} Tilgung bedient; der Rest flie{\ss}t
als Dividende oder Aktienr\"uckkauf an die Eigent\"umer. Der
Entschuldungspfad des vorigen Abschnitts verwendet dieselben Mittel zur
vorzeitigen Tilgung. Beide Rechnungen sind deshalb Alternativen und keine
Erg\"anzungen; dasselbe Geld kann nur einmal ausgegeben werden.}

\annahme{\textbf{Ausstiegsbewertung.} Der Verkaufserl\"os ist das Ergebnis
je Aktie des letzten Jahres, multipliziert mit dem
\Faktor{\FlowVielfaches}-fachen. Dieses Vielfache ist nicht gew\"ahlt,
sondern das heutige: Es ergibt sich aus dem Einstiegskurs von
\USDje{\FlowEinstiegskurs} geteilt durch das verw\"asserte
Pro-forma-Ergebnis je Aktie 2025 von \USDje{\ProFormaEPSdil}.\quelleGB{2025}{64}\quelleMerken{s14gb}
Die Bewertung bleibt im Basisfall also unver\"andert; der Anleger verdient
nur, was das Unternehmen erwirtschaftet, nicht was der Markt ihm zubilligt.
Tabelle~\ref{tab:flowmultiple} variiert dieses Vielfache.}

\begin{table}[htbp]
\centering
\caption{Zahlungsreihe je Aktie in USD \"uber f\"unf Jahre, mittleres
Szenario.}
\label{tab:flowanleger}
\begin{tabular}{lrl}
\toprule
Jahr & Zahlung & Bemerkung\\
\midrule
\tabellenkoerper{data/flow_anleger}
\bottomrule
\end{tabular}
\end{table}

\FloatBarrier

\begin{table}[htbp]
\centering
\caption{Interner Zinsfu{\ss} in Prozent nach Horizont und
Wachstumsszenario, Ausstiegsvielfaches unver\"andert bei
\num{\FlowVielfaches}.}
\label{tab:flowirr}
\begin{tabular}{lrrr}
\toprule
Horizont & \Pct{\HJOrganisch} (H1 2026) & \Pct{\OrganischWachstum} (2025) & \Pct{\OrganischWachstumVJ} (2024)\\
\midrule
\tabellenkoerper{data/flow_irr}
\bottomrule
\end{tabular}
\end{table}

\begin{table}[htbp]
\centering
\caption{Interner Zinsfu{\ss} in Prozent nach Horizont und
Ausstiegsvielfachem, mittleres Wachstumsszenario.}
\label{tab:flowmultiple}
\begin{tabular}{lrrr}
\toprule
Horizont & \num{\FlowVielfachesUnten}-fach & \num{\FlowVielfaches}-fach & \num{\FlowVielfachesOben}-fach\\
\midrule
\tabellenkoerper{data/flow_irr_multiple}
\bottomrule
\end{tabular}
\end{table}

\FloatBarrier

\subsection{Vergleichsma{\ss}stab und Fazit}

Ein interner Zinsfu{\ss} allein tr\"agt kein Urteil; er braucht eine
Gr\"o{\ss}e, die er schlagen muss. Diese Analyse nimmt daf\"ur die
Eigenkapitalrendite des Unternehmens selbst: Das Konzernergebnis von
\MioUSD{\Konzernergebnis} auf das durchschnittliche Eigenkapital des Jahres
2025 bezogen ergibt \Pct{\EKRenditeDurchschnitt}. Ein Kapitalkostensatz
nach dem Kapitalmarktmodell wird bewusst nicht gebildet: Weder Beta noch
Marktrisikopr\"amie stehen in einer der Prim\"arquellen, und ein
gesch\"atzter Satz w\"are genau die plausibel aussehende Erfindung, die
diese Analyse vermeidet.

\bk{Gemessen daran ist das Ergebnis unbequem knapp. Im mittleren Szenario
liegt der interne Zinsfu{\ss} bei \Pct{\FlowIRRBFuenf} \"uber f\"unf,
\Pct{\FlowIRRBZehn} \"uber zehn und \Pct{\FlowIRRBFuenfzehn} \"uber f\"unfzehn Jahre.
Der Ma{\ss}stab von \Pct{\EKRenditeDurchschnitt} wird damit auf f\"unf Jahre
knapp \"ubertroffen und auf zehn und f\"unfzehn Jahre knapp verfehlt. Der
heutige Kurs preist das Unternehmen also ungef\"ahr so, dass ein Anleger
genau die Rendite erh\"alt, die das Unternehmen auf sein eigenes Kapital
erwirtschaftet, und keinen Aufschlag darauf.

Anders als beim Verschuldungsgrad ist die Rangfolge hier nicht stabil,
sondern kippt an der Wachstumsannahme. Bleibt es beim negativen organischen
Wachstum des ersten Halbjahres 2026, sinkt der Zinsfu{\ss} auf
\Pct{\FlowIRRAFuenf} bis \Pct{\FlowIRRAFuenfzehn} und liegt \"uber jeden Horizont
deutlich unter dem Ma{\ss}stab. Kehrt das Wachstum auf die Rate von 2024
zur\"uck, steigt er auf \Pct{\FlowIRRCFuenf} bis \Pct{\FlowIRRCFuenfzehn}. Zwischen
dem schlechtesten und dem besten der drei vom Unternehmen selbst
berichteten Wachstumsjahre liegen \"uber f\"unf Jahre rund sechzehn
Prozentpunkte Rendite. Der Fall entscheidet sich damit an einer einzigen
Gr\"o{\ss}e, und diese Gr\"o{\ss}e ist zuletzt negativ geworden.

Das Ausstiegsvielfache wirkt in dieselbe Richtung, aber schw\"acher und mit
zunehmendem Horizont abnehmend: \"Uber f\"unf Jahre reicht die Spanne von
\Pct{\FlowIRRuntenFuenf} bis \Pct{\FlowIRRobenFuenf}, \"uber f\"unfzehn nur noch von
\Pct{\FlowIRRuntenFuenfzehn} bis \Pct{\FlowIRRobenFuenfzehn}. Je l\"anger die Haltedauer,
desto weniger h\"angt das Ergebnis davon ab, was der Markt am Ende zahlt,
und desto mehr davon, was das Unternehmen verdient. Das ist der einzige
Befund dieses Teils, der in beide Richtungen beruhigt.}
